\section{Stable PCRs and exact algebraic signs}
\label{app:technical}

\subsection{Stable rotation cycles}

\begin{lemma}[Fixed rotations in a stable PCR]
\label{lem:stable-pcr}
Let $C$ be an RC-stable PCR at even length $k$, over any fixed-point-free
complemented alphabet, and let $d$ be its minimal period.  Then $d$ is even
and $C$ contains exactly two reverse-complement-fixed words.
\end{lemma}

\begin{proof}
Choose $w\in C$ and $s\pmod d$ with $\RC(w)=R^sw$.  In cyclic position
notation,
\[
  w_r=\overline{w_{s-1-r}}.
\]
The positional involution $r\mapsto s-1-r$ cannot have a fixed point.  If
$d$ were odd, $2$ would be invertible modulo $d$ and
$2r\equiv s-1\pmod d$ would have a solution, a contradiction.  Thus $d$ is
even.  For even $d$, that congruence is solvable exactly when $s$ is odd, so
$s$ is even.

A rotation $R^iw$ is reverse-complement fixed precisely when
\[
  \RC(R^iw)=R^{-i}\RC(w)=R^{s-i}w=R^iw,
\]
or $2i\equiv s\pmod d$.  Since $d$ and $s$ are even, this congruence has
exactly $\gcd(2,d)=2$ solutions modulo $d$.
\end{proof}

\begin{lemma}
On words of odd length, reverse complement has no fixed point.
\end{lemma}

\begin{proof}
The central letter of a fixed odd-length word would equal its complement.
\end{proof}

These facts permit equivariant representatives on stable zero PCRs, but do
not by themselves imply global acyclicity; Theorem~\ref{thm:arbitrary-zero}
is still essential.

\subsection{Polynomial-bit exact sign determination}

We justify Corollary~\ref{cor:complexity}.  In this subsection only,
specialize the real weight map used by the existence proof to an
integer-valued map.  Put
\[
  c=2\cos(2\pi/k),\qquad d=\varphi(k)/2,
  \qquad B=\max_{a\in\A}|s(a)|.
\]
For $k\geq4$, $c$ is an algebraic integer of degree $d$.  If $\Psi_k$ is its
monic minimal polynomial, it is recovered from the reciprocal cyclotomic
identity
\begin{equation}
  \Phi_k(z)=z^d\Psi_k(z+z^{-1}).
  \label{eq:reciprocal-cyclotomic}
\end{equation}
Classical minimal-polynomial formulas provide the same representation
\cite{WatkinsZeitlin1993}.
Among the real conjugates of $c$, the chosen value is the largest one; this
property gives a canonical exact isolating interval for the root of
$\Psi_k$ at which all subsequent signs are evaluated.

Define integer polynomials
\[
  C_0=2,\ C_1=x,\ C_{j+1}=xC_j-C_{j-1},
  \qquad
  S_0=0,\ S_1=1,\ S_{j+1}=xS_j-S_{j-1}.
\]
At $x=c$,
\[
  C_j(c)=2\cos(2\pi j/k),
  \qquad
  S_j(c)=\frac{\sin(2\pi j/k)}{\sin(2\pi/k)}.
\]
Since $\sin(2\pi/k)>0$, every sign or zero used in the seed and midpoint
rules reduces to an integer polynomial at $c$:
\begin{align*}
  2\operatorname{Re}P_s(w)
    &=\sum_j s(w_j)C_{j+1}(c),\\
  \operatorname{sign}\operatorname{Im}P_s(w)
    &=\operatorname{sign}\sum_j s(w_j)S_{j+1}(c),\\
  \operatorname{sign}\lambda_s(u)
    &=\operatorname{sign}\sum_j s(u_j)S_{j+1}(c),\\
  2r_s(u)
    &=\sum_j s(u_j)\bigl(C_{j+1}(c)-2\bigr).
\end{align*}
Reducing a polynomial modulo $\Psi_k$ decides exact zero: the value vanishes
at $c$ precisely when the remainder is the zero polynomial.

For a nonzero remainder $g$, each conjugate $c_i$ of $c$ lies in $[-2,2]$.
The elementary bounds
$|C_j(c_i)|\leq2$ and $|S_j(c_i)|\leq j$ imply the common safe estimate
\[
  |g(c_i)|\leq M:=2Bk(k+1).
\]
Because $g(c)$ is a nonzero algebraic integer, its field norm is a nonzero
rational integer.  Therefore
\begin{equation}
  |g(c)|\geq M^{-(d-1)}.
  \label{eq:norm-gap}
\end{equation}
Certified interval evaluation below half this gap determines the sign.

The coefficients of $\Phi_k$, $\Psi_k$, the recurrences, and their remainders
have bit length polynomial in $k$ and $\log(B+1)$.  Standard
Sturm--Habicht real-root isolation and sign evaluation operate in polynomial
bit complexity \cite{EmirisTsigaridas2006}.  Equation~\eqref{eq:norm-gap}
requires only $O(k(\log B+\log k))$ refinement bits.  Scanning at most $k$
rotations handles a deterministic zero-PCR tie.  This proves the claimed
polynomial preprocessing, query, and workspace bounds without a real-RAM or
floating-point assumption.
